\scanpage{12}
\mspage{\unsure{6}}
Programme du \deleted{paradigme \unsure{de la théorie des|arithmétique des}} \inserted{``\unsure{gruppentheoretisches}} Gleichnis'' (paradigme en termes de théorie des groupes) de l'\inserted{arithmétique :}\deleted{\unsure{arithm.} géométrie algébrique absolue (en car. 0).}

a) Donner une description ``groupo-théorétique'' complète de $\Gamma$.

b) Donner une description complète de tous les éléments de structure énumérés ci-dessus, dans des cas typiques importants pour $(X,D,J)$, dont nous indiquons quelques-uns plus bas.

\margin{On \unsure{tâchera} de prendre $X$ \unsure{propre autant} que possible (quitte à \unsure{compléter}\ldots)}

\uline{NB} Quand on a les éléments de structure envisagés pour $(X,D,J)$, on les a aussi automatiquement pour $\forall$ rev.\inserted{étale} fini de $X'\ldots$ Il faudrait expliciter ce point avec soin, et expliciter également dans les $\unsure{=|2}$ optiques les foncteurs ``bouchage de trous'' (en rajoutant \inserted{certains} des $D_i$ qu'on avait \unsure{enlevés|exclus}) ou ``forage des trous'' \inserted{si $\dim X=1$} (en enlevant certains des $\xi\in J$).
\mspage{11}
